\section{System Design}
\label{sec:system-design}

This section presents the system design of the AI Football Prediction System. The system is designed as a web-based machine learning application that predicts whether an upcoming football match is likely to have over or under 2.5 total goals. The application consists of a React web frontend, a FastAPI backend, a data/model pipeline, local machine learning artifacts, and an analytics database.

The main design objective is to separate the online prediction application from the offline machine learning workflow. The web application is responsible for user interaction, while the backend API is responsible for business logic, data access, analytics tracking, and prediction execution. Model training is handled as an offline process, and the trained model files are later loaded by the backend when users request predictions.

\subsection{High-Level Architecture}
\label{subsec:high-level-architecture}

The system follows a layered architecture. The frontend layer contains the React web application. The backend layer contains the FastAPI service. The data layer contains CSV/JSON football data, trained model files, and an analytics database. The machine learning pipeline is executed separately to prepare the data and model artifacts used by the backend.

\begin{figure}[htbp]
    \centering
    % \includegraphics[width=0.95\textwidth]{figures/system_architecture.png}
    \caption{High-level architecture of the web-based football prediction system.}
    \label{fig:system-architecture}
\end{figure}

The web frontend communicates with the backend through REST API calls. The backend reads football data from local files under the \texttt{data} directory and loads trained models from the \texttt{models} directory. For user accounts and activity analytics, the backend supports SQLite for local execution and PostgreSQL when the system is executed with Docker Compose.

\subsection{Main Components}
\label{subsec:main-components}

\subsubsection{Web Frontend}

The web frontend is implemented with React and Vite and is located in \texttt{apps/web}. It provides role-based screens for users, data scientists, and administrators. Normal users can browse leagues, teams, upcoming matches, team details, and prediction results. Data scientists can view model-related dashboard information. Administrators can view user activity analytics such as frequently accessed teams and players.

The frontend does not perform machine learning inference locally. Instead, it sends HTTP requests to the backend API and renders the returned data. The API base URL is configurable through \texttt{VITE\_API\_BASE\_URL}; in local development it points to:

\begin{verbatim}
http://127.0.0.1:8000/api/v1
\end{verbatim}

\subsubsection{Backend API}

The backend is implemented with FastAPI and is located in \texttt{apps/api}. It is the central runtime component of the system. Its responsibilities include:

\begin{itemize}
    \item exposing REST endpoints for the web frontend;
    \item loading football league, team, match, and player-related data;
    \item loading trained machine learning models from disk;
    \item preparing feature vectors for upcoming matches;
    \item returning Over/Under 2.5 prediction results;
    \item storing user accounts and activity analytics;
    \item enforcing simple role-based access control.
\end{itemize}

The backend exposes endpoints under the \texttt{/api/v1} prefix. The main route groups are:

\begin{itemize}
    \item \texttt{/api/v1/user/...}: user-facing football browsing and prediction endpoints;
    \item \texttt{/api/v1/data-scientist/...}: model dashboard and analysis endpoints;
    \item \texttt{/api/v1/admin/...}: administrative analytics endpoints.
\end{itemize}

The backend also exposes \texttt{/health}, which is used by Docker Compose to check whether the API service is running.

\subsubsection{Data Repository}

The football dataset is stored as local files. Historical raw data is stored in \texttt{data/raw}, while preprocessed model-ready data is stored in \texttt{data/processed}. Upcoming matches are stored in \texttt{data/next\_matches.json}. This file-based design keeps the data pipeline simple and makes the system easier to reproduce in a local or academic environment.

The main data artifacts are:

\begin{itemize}
    \item \texttt{data/raw/<league>\_merged.csv}: merged historical match data;
    \item \texttt{data/processed/<league>\_merged\_preprocessed.csv}: processed feature data;
    \item \texttt{data/next\_matches.json}: upcoming match data used by the web application;
    \item \texttt{data/final\_predictions.txt}: optional text output generated by the prediction script.
\end{itemize}

\subsubsection{Model Layer}

The trained models are stored as pickle files in the \texttt{models} directory. For each league, the system can store a Random Forest model and an ensemble voting classifier. The backend prediction service loads the requested model and uses processed football statistics to prepare the input feature vector for prediction.

The expected model artifacts are:

\begin{itemize}
    \item \texttt{models/<league>\_random\_forest.pkl};
    \item \texttt{models/<league>\_voting\_classifier.pkl}.
\end{itemize}

The backend currently supports two model types:

\begin{itemize}
    \item \texttt{random\_forest};
    \item \texttt{ensemble}.
\end{itemize}

\subsubsection{Analytics Database}

The system stores user accounts and analytics events in a relational database. In local development, the backend falls back to SQLite at \texttt{apps/api/app/db/analytics.db}. When the application is run with Docker Compose, the backend uses PostgreSQL through the \texttt{DATABASE\_URL} environment variable.

The analytics database stores:

\begin{itemize}
    \item registered users;
    \item team view events;
    \item player view events;
    \item match view events.
\end{itemize}

This design allows the project to remain lightweight for local development while still supporting a more production-like database setup in Docker.

\subsection{Data and Model Pipeline}
\label{subsec:data-model-pipeline}

The machine learning workflow is executed offline before running the prediction application. This pipeline prepares the input data, trains the models, and stores the generated artifacts for backend usage.

\begin{figure}[htbp]
    \centering
    % \includegraphics[width=0.95\textwidth]{figures/data_pipeline.png}
    \caption{Offline data and model training pipeline.}
    \label{fig:data-pipeline}
\end{figure}

The pipeline contains three main stages:

\begin{enumerate}
    \item \textbf{Data acquisition}: \texttt{scripts/data\_acquisition.py} downloads historical football CSV files and merges multiple seasons for each league.
    \item \textbf{Data preprocessing}: \texttt{scripts/data\_preprocessing.py} cleans data, engineers football statistics, handles missing values, and selects useful features.
    \item \textbf{Model training}: \texttt{scripts/train\_models.py} trains Random Forest and ensemble models, then saves them as pickle files in \texttt{models}.
\end{enumerate}

The helper script \texttt{run.sh} wraps the pipeline commands. A full pipeline run can be executed with:

\begin{verbatim}
./run.sh pipeline
\end{verbatim}

After the pipeline has been executed once, the application can reuse the generated \texttt{data} and \texttt{models} artifacts. The pipeline only needs to be rerun when the dataset, features, or model training logic changes.

\subsection{Prediction Request Flow}
\label{subsec:prediction-request-flow}

When a user requests a match prediction, the web application sends a request to the FastAPI backend. The backend resolves the selected match, records the user activity, prepares the model input, loads the trained model, and returns the prediction result.

\begin{figure}[htbp]
    \centering
    % \includegraphics[width=0.9\textwidth]{figures/prediction_sequence.png}
    \caption{Sequence diagram for a match prediction request.}
    \label{fig:prediction-sequence}
\end{figure}

The prediction flow is:

\begin{enumerate}
    \item The user selects an upcoming match in the web application.
    \item The frontend calls the backend prediction endpoint.
    \item The backend finds the selected match from the upcoming match data.
    \item The backend records a match view event in the analytics database.
    \item The prediction service loads processed team statistics.
    \item The prediction service builds a numerical feature vector.
    \item The trained model predicts Over/Under 2.5 and probability values.
    \item The backend returns the result to the frontend.
    \item The frontend displays the prediction to the user.
\end{enumerate}

\subsection{Use Case Design}
\label{subsec:use-case-design}

The system supports three roles: User, Data Scientist, and Admin. These roles access the same web application but use different features.

\begin{figure}[htbp]
    \centering
    % \includegraphics[width=0.9\textwidth]{figures/use_case_diagram.png}
    \caption{Use case diagram of the football prediction system.}
    \label{fig:use-case-diagram}
\end{figure}

The main use cases are:

\begin{itemize}
    \item \textbf{Login/Register}: create an account or log in.
    \item \textbf{View leagues}: browse supported leagues.
    \item \textbf{View teams}: view teams in a selected league.
    \item \textbf{View upcoming matches}: inspect upcoming fixtures.
    \item \textbf{View team details}: inspect recent and upcoming matches for a team.
    \item \textbf{View prediction}: view Over/Under 2.5 prediction for a selected match.
    \item \textbf{View model dashboard}: inspect model-related metrics and summaries.
    \item \textbf{View analytics}: inspect top viewed teams, players, matches, and users.
\end{itemize}

In the use case diagram, the \textit{User} actor should be connected to browsing leagues, teams, upcoming matches, team details, and match predictions. The \textit{Data Scientist} actor should be connected to the model dashboard. The \textit{Admin} actor should be connected to analytics views. All roles should be connected to login/register.

\subsection{Deployment Design}
\label{subsec:deployment-design}

The system can run in two modes. In local development mode, the backend uses SQLite and the web/API services are started separately through \texttt{run.sh}. In Docker mode, Docker Compose starts the React web server, FastAPI backend, and PostgreSQL database together.

\begin{figure}[htbp]
    \centering
    % \includegraphics[width=0.9\textwidth]{figures/deployment_diagram.png}
    \caption{Docker Compose deployment design.}
    \label{fig:deployment-diagram}
\end{figure}

The Docker deployment contains:

\begin{itemize}
    \item \textbf{web}: serves the built React application through Nginx;
    \item \textbf{api}: runs the FastAPI backend with Uvicorn;
    \item \textbf{postgres}: stores users and analytics data.
\end{itemize}

The \texttt{data} and \texttt{models} directories are mounted into the API container. This means the source code can be pushed to GitHub without committing large generated CSV/model files. After cloning the project, the pipeline can be rerun to regenerate these artifacts.

\subsection{Figures to Prepare}
\label{subsec:figures-to-prepare}

The following figures should be included in the report:

\begin{enumerate}
    \item \textbf{System architecture diagram} (\texttt{figures/system\_architecture.png}): show the React Web App calling the FastAPI Backend. Behind the backend, show File Data Storage, Model Files, and Analytics DB. Also show the offline ML pipeline producing data/model artifacts.
    \item \textbf{Data pipeline diagram} (\texttt{figures/data\_pipeline.png}): show Historical CSV Download, Raw Data, Preprocessing, Feature Engineering/Selection, Model Training, and Saved Model Artifacts.
    \item \textbf{Prediction sequence diagram} (\texttt{figures/prediction\_sequence.png}): show User, Web Frontend, FastAPI Backend, Data Repository, Prediction Service, ML Model, and Analytics DB.
    \item \textbf{Use case diagram} (\texttt{figures/use\_case\_diagram.png}): show actors User, Data Scientist, and Admin with their corresponding use cases.
    \item \textbf{Deployment diagram} (\texttt{figures/deployment\_diagram.png}): show Docker Compose with Web container, API container, PostgreSQL container, and mounted \texttt{data}/\texttt{models} folders.
\end{enumerate}
